% Gesamttest «Lineare und quadratische Gleichungen» — Quelle des PDFs (python3 scripts/build-lp-pdf.py lineare-quadratische).
% Musterlösung und Raster: bewertungspaket.tex. Alle Werte mit python3 nachgerechnet (06.10.2026).
% Kein \lt/\gt — pdfLaTeX kennt sie nicht, < und > tun es.
% Neu zusammengestellt 06.10.2026 nach der Prüfung (TODO.md H1–H3, M6): Kombinationen und Transfer statt
% Wiederholung der Kapitelmodelle. Bewusst 30 statt «rund 20» Minuten (HOWTO §9): 8 Aufgaben ohne Rechner,
% darunter zwei Parameterdiskussionen und eine Begründungsaufgabe (G6).
% Nachprüfung 06.10.2026 (TODO.md N-H1, N-M2, N-M4): G3, G4, G5, G6(d), G8 weiter weg von den Kapitelmodellen.
\documentclass[11pt]{article}
\usepackage{../lp-druck}
\renewcommand{\lpname}{Lineare und quadratische Gleichungen}
\renewcommand{\lpteilgebiet}{2.2}
\renewcommand{\lpbereich}{Grundlagenfach}
\renewcommand{\lpfuss}{Gesamttest · 25 Punkte}
\begin{document}
\lpkopf{Gesamttest}{%
  \vspace{4pt}\noindent\begin{tabularx}{\linewidth}{@{}X X@{}}
  Code (kein Name): \hrulefill & Punkte: \hrulefill\ / 25
  \end{tabularx}}

\begin{lpinfo}{Anleitung}
\textbf{Zeit:} rund 30 Minuten. \textbf{Hilfsmittel:} keine — kein Taschenrechner (RLP 2.2: «auch ohne Hilfsmittel»).
Schreib zu jeder Aufgabe den \textbf{Rechenweg} auf — bewertet wird der Weg. Gib Lösungen als Lösungsmenge $\mathbb{L}$ an; trenne die Elemente mit Strichpunkt, z.\,B. $\{-2;\, 3\}$.
Danach: Lösung scannen oder fotografieren (ohne Namen und Standort) und mit dem \emph{Bewertungspaket} bewerten lassen.
\end{lpinfo}

\lpteil{Teil A · Umformen und Nullprodukt}{Kapitel 1, 2 und 5 · 8 Punkte}

\begin{lpaufgabe}{G1}{Lineare Gleichung mit Brüchen}{3 P}
Löse und mach die Probe: $\;\dfrac{x + 2}{4} - (x - 3) = \dfrac{x - 1}{2}$
\lpplatz{7}
\end{lpaufgabe}

\begin{lpaufgabe}{G2}{Lösungsfälle}{2 P}
Gegeben ist $3(x - 2) = 3x + c$. Für welches $c$ ist jede Zahl eine Lösung, für welche $c$ gibt es keine Lösung?
Begründe mit der Aussage, die beim Umformen übrig bleibt.
\lpplatz{4}
\end{lpaufgabe}

\begin{lpaufgabe}{G3}{Nullprodukt und Ausklammern}{3 P}
\begin{enumerate}[label=(\alph*),leftmargin=*,itemsep=1pt]
\item Löse $(2x - 1)(x + 3) = 0$.
\item Jemand löst $3x^2 = 7x$ so: «beidseitig durch $x$ teilen: $3x = 7$, also $x = \tfrac{7}{3}$». Welche Lösung geht dabei verloren, und warum? Löse die Gleichung richtig.
\end{enumerate}
\lpplatz{5}
\end{lpaufgabe}

\lpteil{Teil B · Quadratische Gleichungen lösen}{Kapitel 3 und 4 · 11 Punkte}

\begin{lpaufgabe}{G4}{Quadratische Ergänzung mit Zahl \texorpdfstring{$c$}{c}}{3 P}
Gegeben ist $x^2 - 10x + c = 0$.
\begin{enumerate}[label=(\alph*),leftmargin=*,itemsep=1pt]
\item Bring die Gleichung mit quadratischer Ergänzung auf die Form $(x - 5)^2 = \ldots$
\item Lies daraus ab: Für welche $c$ gibt es zwei, eine oder keine Lösung?
\item Löse die Gleichung für $c = 16$.
\end{enumerate}
\lpplatz{5}
\end{lpaufgabe}

\begin{lpaufgabe}{G5}{Mitternachtsformel, exakt}{3 P}
Löse $x\,(x - 2) = 1$ mit der Mitternachtsformel. Gib $D$ an und die Lösungen exakt (mit Wurzel).
\lpplatz{5}
\end{lpaufgabe}

\begin{lpaufgabe}{G6}{Verfahren begründen und prüfen}{5 P}
Nenne in (a) bis (c) für jede Gleichung das schnellste Verfahren und begründe die Wahl in einem Satz. Nicht lösen.
\begin{enumerate}[label=(\alph*),leftmargin=*,itemsep=1pt]
\item $4x^2 - 9 = 0$\par\vspace{1.4\baselineskip}\noindent\rule{\linewidth}{0.3pt}
\item $x^2 + 6x + 9 = 0$\par\vspace{1.4\baselineskip}\noindent\rule{\linewidth}{0.3pt}
\item $2x^2 + 5x - 1 = 0$\par\vspace{1.4\baselineskip}\noindent\rule{\linewidth}{0.3pt}
\item Jemand gibt für $x^2 - 4x - 12 = 0$ die Lösungen $2$ und $6$ an. Prüfe beide mit der Probe und korrigiere, was falsch ist.
\end{enumerate}
\lpplatz{8}
\end{lpaufgabe}

\lpteil{Teil C · Parameterdiskussion}{Kapitel 5 · 6 Punkte}

\begin{lpaufgabe}{G7}{Lineare Gleichung mit Parameter}{3 P}
Löse $k(k - 1) \cdot x = k$ nach $x$ auf ($k \in \mathbb{R}$). Unterscheide alle Fälle und gib jeweils die Lösungsmenge an.
\lpplatz{9}
\end{lpaufgabe}

\begin{lpaufgabe}{G8}{Quadratische Gleichung mit Parameter}{3 P}
Untersuche $(k - 1)\,x^2 - 4x + 2 = 0$ ($k \in \mathbb{R}$): Für welche $k$ gibt es zwei, eine oder keine Lösung? Gib die Lösung an, wo es genau eine gibt.
\lpplatz{9}
\end{lpaufgabe}
\end{document}
